\section{Representative Systems and Design Patterns}

Table~\ref{tab:systems} positions representative systems in the proposed taxonomy.

\begin{table}[t]
\caption{Representative GPU-integrated I/O systems.}
\label{tab:systems}
\centering
\begin{tabular}{@{}p{2.3cm}p{2.2cm}p{2.3cm}p{5.4cm}@{}}
\toprule
System & Control Plane & Data Plane & Key Architectural Characteristic \\
\midrule
GPUfs \cite{silberstein2013gpufs} & CPU + GPU runtime & Staged / managed & File-like interface in GPU programs; emphasizes programmability. \\
GDS \cite{nvidia_gds_overview} & CPU-coordinated & Direct DMA to GPU & Reduces host bounce buffering; production integration with CUDA stack. \\
BaM \cite{qureshi2023bam} & GPU-initiated + CPU setup & Direct and demand-driven & Fine-grain, high-throughput GPU-initiated storage access model. \\
GeminiFS \cite{geminifs2025} & Hybrid distributed control & GPU-first data path & Optimizes distributed DL checkpoint and ingest behavior. \\
\botrule
\end{tabular}
\end{table}

Across these systems, three recurring design patterns emerge. \textbf{Pattern P1: split-phase control and transfer}. Initialization and protection remain host-resident, while transfer pipelines are aggressively offloaded. \textbf{Pattern P2: queue-depth amplification}. Throughput scaling relies on many in-flight requests, often constrained by completion handling strategy rather than submission bandwidth. \textbf{Pattern P3: format-aware batching}. Effective systems co-design storage layout and GPU kernel consumption order, reducing random-access amplification.

Notably, direct paths are not universally superior. Under metadata-heavy workloads, host-side namespace checks and synchronization can dominate. Under low queue depth, PCIe/NVMe headroom is underutilized; under extreme depth, tail latency rises and harms synchronous kernel phases \cite{nvme_spec,nvidia_gds_best_practices}.
